\section{Introduction}
\label{sec:intro}

Street-view perception models can predict subjective attributes such as
safety, wealth, or liveliness at scale, enabling geographic mapping and
correlational analysis at a scope previously unavailable to urban
researchers~\cite{salesses2013,naik2014,dubey2016}. Yet these models
remain limited as explanatory tools: they do not answer the question
central to urban planning and design, \emph{for this specific scene,
which visual elements are responsible for the perception, and what
targeted modification would plausibly shift it?}

Existing explainability approaches such as saliency maps, SHAP, and
concept probes remain largely correlational: they identify pixels or
features associated with a score without validating that manipulating
those features changes human judgement~\cite{adebayo2018,kim2018}. The
distinction matters. A correlational explanation may point to a visually
dominant cue that is not causally active; an interventional test instead
asks whether a proposed localized change actually shifts perception.

This gap is particularly consequential in urban design. Empirical work
using street-view imagery finds that greenery, street-facing windows,
and maintenance cues contribute positively to perceived
safety~\cite{li2015greenery,denadai2016lively}, while intervention
studies and lighting research emphasize the role of disorder reduction
and visibility~\cite{jiang2018brokenwindows,portnov2020lighting}. Yet
these studies do not isolate the contribution of individual cues through
valid counterfactuals. More fundamentally, following Jacobs and
defensible-space theory~\cite{jacobs1961,newman1972}, perceived safety
is often understood as an emergent property of multiple visible cues
acting simultaneously. This raises a question prior work has not
directly tested: does a given scene admit one dominant visual pathway to
improved safety perception, or multiple substitutable levers that could
each, in isolation, plausibly shift perception?

We propose an \emph{interventional attribution} protocol that makes
these questions empirically testable. Rather than treating arbitrary
image regions as the explanatory unit, we define structured lever
interventions grounded in a curated ontology of manipulable urban cues,
each specifying a semantic concept, spatial support, intervention
direction, and constrained edit template. Candidate interventions are
generated via prompt-conditioned diffusion editing and filtered by a
vision-language critic, retaining only edits that appear localized,
plausible, and same-place preserving. This yields a controlled
procedure for testing which urban cues can be validly realized for a
specific scene and whether scenes admit multiple substitutable candidate
levers.

Our empirical questions are therefore twofold:

\begin{quote}
\textit{Which localized urban cues can be validly realized under
prompt-conditioned diffusion editing, and which of them directionally
improve an auxiliary safety proxy? Across scenes, how many valid
substitutable levers typically exist: is safety perception visually
overdetermined, or concentrated in a smaller set of intervention
pathways?}
\end{quote}

Generative editors are an imperfect intervention mechanism: they may
fail to implement the intended change or introduce correlated,
non-target modifications that also affect perception. We therefore treat
generated edits as hypotheses to be audited rather than as faithful
evidence by default. To make this tractable, we also restrict the
planner to a curated ontology of manipulable urban cues informed by the
urban-perception literature. Without such a scaffold, free-form
intervention proposals become unbounded and difficult to compare across
scenes. The ontology thus serves not only as a substantive prior, but
also as an experimental design device that converts open-ended
generation into a controlled, auditable search problem.

Our contributions are threefold:
\begin{enumerate}
    \item We introduce a \textbf{lever-based interventional attribution}
    framework for street-view perception, in which the explanatory atom
    is a structured lever intervention
    $l=(c,s,d,\tau)$ rather than an arbitrary image region.

    \item We propose a \textbf{scene-specific candidate construction and
    evaluation pipeline} that instantiates lever interventions from a
    curated ontology of manipulable urban cues, realizes them through
    constrained prompt-only edits, and audits them for same-place
    preservation, locality, realism, and plausibility.

    \item We provide a \textbf{reproducible pilot empirical study} on
    urban safety perception on a five-city subset, characterizing which
    urban cue families are currently realizable and which produce the
    strongest directional shifts in an auxiliary perception score, while
    positioning human pairwise evaluation as the future endpoint.
\end{enumerate}
